\chapter{关键技术挑战与系统设计}

\section{如何把学生设计转化为可执行任务}

操作系统 Lab 通常只给出功能目标，例如增加系统调用、实现页分配器或接入块设备。这样的描述适合学生阅读，但平台还需要确定当前任务涉及哪些模块、允许修改哪些文件、必须保持哪些性质，以及用什么结果判断完成。若这些信息只存在于问答记录中，Agent、测试工具和教师会对同一个任务产生不同理解。

VeriSpecOSLab 用五类规格分别记录这些决定。DesignSpec 说明体系结构和内核组织。ModuleSpec 说明当前模块的职责、可修改范围、应保持的性质和验收项目。InterfaceSpec 说明系统调用、驱动等跨模块接口。GoalSpec 记录可以度量的扩展目标。SpecPatch 记录确需跨模块修改时的原因和影响范围。学生在讨论设计后亲自填写规格，平台负责检查结构、引用和测试对应关系。

实现中的 \code{buildNormalizedSpecBundle} 读取全部规格文件，校验 YAML 与 schema，统一 Spec ID 和文件路径，解析模块、接口、操作与测试要求，并为每个输入文件计算摘要。规格存在错误时，平台在 Agent 开始编码前停止任务。解析通过后，当前 ModuleSpec 的文件范围、性质和测试要求成为实现、验证和报告共同使用的输入。学生的设计由此转化为一份可以被程序检查的 Lab 任务。

\section{如何保证 AI 只完成当前 Lab}

在提示词中要求模型``只修改当前模块''无法形成可靠边界。模型可能为了修复编译错误改动其他模块，也可能顺手实现课程后续内容。平台首先要求当前规格已经提交，且学生工作目录没有未提交修改。随后从目标 ModuleSpec 的 \code{owns} 计算允许修改的文件。确需跨模块调整时，学生必须先提交 SpecPatch，平台再把其中声明的影响范围加入本次任务。

实现助手运行在从当前 Git 提交创建、与学生目录分开的临时 Git 工作目录中。工具配置只向模型提供允许操作的路径，任务提示同时写明当前模块和测试要求。模型运行结束后，平台使用 Git 读取已跟踪文件的差异和未跟踪文件列表，并逐项检查实际变化是否位于允许范围。这个检查依据工作目录中的真实文件，不使用模型自行报告的修改清单。

平台还会恢复模型对 \code{vos.yaml} 的直接修改。新增测试必须通过结构化结果交给平台，再由平台检查并写入运行配置。本地补充测试的内容同样通过结果字段返回，不能作为源文件留在 Git 工作目录中。只有修改范围、规格解析、项目构建和全部检查均通过后，平台才在临时工作目录创建提交。回到学生项目时，平台再次确认原始提交没有变化、工作目录仍然干净，再把经过验证的提交直接接到学生项目的当前提交之后。失败任务保留补丁与日志，学生当前目录不会混入未通过的代码。

\section{如何判断实现符合学生设计}

一份代码能够编译，只能说明编译器接受了当前实现。它仍可能遗漏规格中的边界情况，使用错误机制，或者提供与当前模块无关的测试。VeriSpecOSLab 因此要求实现助手提交 \code{student\_implementation\_result.v1} 结构化结果。成功结果需要明确标记状态，并同时给出公开测试、接口契约测试、固定随机种子的模糊测试、有限执行轨迹测试和本地补充测试。

平台逐项检查测试标识是否与已有项目冲突，命令引用的文件是否存在，工作目录和产物路径是否合法，最长运行时间是否有效，以及 \code{verifies} 是否引用当前 Spec ID。结构化结果不满足要求时，错误会返回同一个模型会话，模型可以继续修改文件和结果字段。平台不会从自然语言回答中猜测缺失字段，也不会把 \code{failed}、\code{partial} 或 \code{blocked} 当作成功。

测试配置通过检查后，平台将其加入临时工作目录中的运行配置，再重新解析规格与配置。随后执行项目构建和配置中全部检查，并在执行结束后再次检查文件范围。最终状态由结构化结果、Git 实际变化和命令运行结果共同决定。模型给出的文字总结不参与状态转换。这样能够区分``模型认为已经完成''和``当前提交满足规格并通过实际验收''。

\section{如何证明运行结果对应最终提交}

学生可能在未提交文件、旧的构建产物或不同运行配置下得到成功结果。如果报告只保存一段终端输出，教师无法确认这段输出对应哪份代码。平台允许学生在开发过程中随时编译和运行 QEMU，但工作目录存在未提交修改时，结果会被标记为不可提交。正式验证和提交要求当前 Git 提交保持干净。

每次运行都有独立的 Run ID。运行记录保存 Git 提交、规格摘要、运行配置摘要、实际命令、标准输出、错误输出、产物和最终状态。QEMU 结果依据串口中的成功与失败标记判断。开发板结果另外保存板卡、串口、工作负载和代码提交，并将脚本执行状态与教学团队的人工复核状态分开记录。

提交阶段再次读取当前 Git 提交和规格，检查连续哈希审计链，并确认最近一次补充测试验证绑定相同的提交、规格摘要、配置摘要和测试清单摘要。源码通过 \code{git archive HEAD} 提取，因此归档中只包含当前提交的文件。最终压缩包同时保存代码、报告、审计记录、补充测试和提交清单。清单记录各部分的摘要，使评委能够核对运行成功的代码是否就是收到的代码。

\section{如何维护逐步展开的课程案例}

\code{xv6-spec} 使用一条从 Lab 1 到 Lab 10 逐步展开的 Git 历史。每个课程标签都是可以独立检出的累计快照。早期 Lab 只应包含已经讲授的代码、规格、测试和术语。如果后续实现、测试名称或提示文字提前进入早期标签，学生可能在当前任务中直接看到未来答案。

只比较相邻标签的差异无法发现这类问题，因为未来文件可能早已存在，只是暂时没有被使用。课程历史审计因此读取每个标签的完整文件树。脚本检查某类文件最早可以出现在哪个 Lab，搜索尚未讲授的术语，并核对当前阶段应当存在的 Spec ID 和测试 ID。例如，内存管理文件从 Lab 3 出现，进程与系统调用从 Lab 5 出现，开发板文件从 Lab 9 出现。某项内容提前出现或累计内容缺失时，审计立即失败。

这套检查使课程阶段成为可以自动验证的发布条件。课程维护者修改任一标签后，可以重新检查 Lab 1--10 的完整历史。同一案例也由此验证 VeriSpecOSLab 能够在连续课程中处理规格演进、AI 实现、QEMU 测试、开发板运行和最终提交。
